\section{Appendix: Concrete Interfaces}
\label{sec:appendix}

\subsection{OSEK/VDX Checking Options}

The methodology in Section~3 refers to OSEK/VDX-specific model construction without listing
its command-line spellings. The following schematic command shows the documented options for
selecting the OSEK/VDX checking model and the corresponding OIL configuration file.

\begin{lstlisting}[caption={Schematic OSEK/VDX option usage.}]
cbmc main.c --os-api osek --osek-oil app.oil
\end{lstlisting}

\subsection{Example \texttt{goto-cc} Invocation}

The following command illustrates the public project-mode flow on the
\texttt{i8xx-tco-1} benchmark tree. It records the project root, names the interruption
source file explicitly, and emits an interleaving manifest.

\begin{lstlisting}[caption={Example project-mode interleaving analysis command.}]
goto-cc \
  check-src/benchmark-sources/intabs/cases/i8xx-tco-1/improved-pipeline/main.c \
  check-src/benchmark-sources/intabs/cases/i8xx-tco-1/improved-pipeline/isr_define/isr.c \
  --interleaving-project-root \
    check-src/benchmark-sources/intabs/cases/i8xx-tco-1/improved-pipeline \
  --interleaving-source-files \
    check-src/benchmark-sources/intabs/cases/i8xx-tco-1/improved-pipeline/isr_define/isr.c \
  --interleaving-output \
    check-src/benchmark-sources/intabs/cases/i8xx-tco-1/interleaving_pipeline.json \
  -o check-src/benchmark-sources/intabs/cases/i8xx-tco-1/i8xx-tco-1.out
\end{lstlisting}

\subsection{Example \texttt{aib} Invocation}

\begin{lstlisting}[caption={Example AIB rewrite command.}]
aib \
  check-src/benchmark-sources/intabs/cases/i8xx-tco-1/improved-pipeline \
  check-src/benchmark-sources/intabs/cases/i8xx-tco-1/interleaving_pipeline.json \
  out/i8xx-tco-1_injected \
  out/i8xx-tco-1_injected/interleaving_pipeline_effective.json
\end{lstlisting}

\subsection{Illustrative Guarded Call}

The restricted checking model is created by inserting guarded interruption calls only at the
selected program points. The example below is illustrative rather than benchmark-specific.

\begin{lstlisting}[caption={Illustrative guarded interruption call.}]
if (nondet_bool()) isr1(0);
\end{lstlisting}

\subsection{Documented Selection-Artifact Fields}

\begin{table}[H]
  \centering
  \small
  \begin{tabularx}{\linewidth}{@{}lX@{}}
    \toprule
    Field & Meaning \\
    \midrule
    \texttt{project.project\_root} & Root directory used for path scoping and
    report-path relativization. \\
    \texttt{project.interleaving\_source\_files} & Source files whose functions are treated
    as interleaving functions. \\
    \texttt{project.translation\_units} & Effective build inputs associated with the initial
    checking model that produced the selection artifact. \\
    \texttt{interleaving[*].name} & Interleaving function identifier recorded for rewriting. \\
    \texttt{interleaving[*].write\_global\_var} & Global symbols that the interleaving
    function may write. \\
    \texttt{interleaving[*].line\_added\_block\_with\_file} & Project-relative file and line
    pairs where guarded calls may be inserted. \\
    \bottomrule
  \end{tabularx}
  \caption{Documented fields of the selection artifact used between analysis and rewriting.}
  \label{tab:manifest-schema}
\end{table}

\subsection{Manifest Excerpt}

The manifest below is a schematic excerpt showing only the documented public fields. It
illustrates how project metadata and interleaving candidates are separated.

\begin{lstlisting}[caption={Schematic manifest excerpt.}]
{
  "project": {
    "project_root": "/path/to/check-src/benchmark-sources/intabs/cases/i8xx-tco-1/improved-pipeline",
    "interleaving_source_files": ["isr_define/isr.c"],
    "translation_units": [
      "main.c",
      "isr_define/isr.c"
    ]
  },
  "interleaving": [
    {
      "name": "isr1",
      "write_global_var": ["<global_1>", "<global_2>"],
      "line_added_block_with_file": [
        { "file": "main.c", "line": [<line_a>, <line_b>] }
      ]
    }
  ]
}
\end{lstlisting}

This schema is the compatibility boundary between \texttt{interleaving-analysis} and
\texttt{aib}; any field change must be implemented on both sides.
